\documentclass{article}
\usepackage[utf8]{inputenc}
\usepackage[T1]{fontenc}
\usepackage{geometry}
\geometry{
    a4paper,
    left=2.5cm,
    right=2.5cm,
    top=2.5cm,
    bottom=2.5cm
}
\usepackage{graphicx}
\usepackage{float}
\usepackage{hyperref}
\usepackage[english]{babel}
\usepackage{amsmath}
\usepackage{amssymb}
\usepackage{array}
\usepackage{booktabs}
\usepackage{caption}
\usepackage{subcaption}
\usepackage{enumitem}
\usepackage{fancyhdr}
\usepackage{titlesec}
\usepackage{url}
\usepackage{textcomp}
\usepackage{xcolor}
\usepackage{listings}
\usepackage{multirow}
\usepackage{longtable}
\usepackage{tabularx}
\usepackage{times}

% Set up page style
\pagestyle{fancy}
\fancyhf{}
\rfoot{\thepage}
\renewcommand{\headrulewidth}{0pt}

% Title formatting
\titleformat{\chapter}[display]{\normalfont\bfseries}{\chaptertitle\thechapter}{0pt}{\Huge}
\titlespacing*{\chapter}{0pt}{50pt}{40pt}

% Document info
\title{Template Document}
\author{Academic Support Agent}
\date{\today}

\begin{document}
\maketitle

\pagenumbering{roman} %comienza con n mero romano para las p ginas
\maketitle %agrega la portada al documento

%se incluye lo contenido en resumen.tex
%% Content from resumen would be included here


%se incluye lo contenido en abstract.tex
%% Content from abstract would be included here


%se incluye lo contenido en agradecimientos.tex
%% Content from agradecimientos would be included here


%Estas 3 listas se crear n autom ticamente, seg n los cap tulos y secciones del documento 




%incluye lo contenido en nomenclatura.tex
% % Content from nomenclatura would be included here


%incluye lo contenido en C1problema.tex
% Content from C1problema would be included here


%incluye lo contenido en C2marcoteorico.tex
% Content from C2marcoteorico would be included here


%incluye lo contenido en C3objetivos.tex
% Content from C3objetivos would be included here


%incluye lo contenido en C4conclusiones.tex
% Content from C4medologia would be included here


%incluye lo contenido en C5cartagantt.tex
% Content from C5cartagantt would be included here


%Generacion automatica de la lista de referencias 
 %estilo APA
% %estilo tipo EEE 
 %indica nombre del archivo de referencias
\end{document}